\documentclass[tikz,border=6pt]{standalone}
\usepackage{ctex}
\usepackage{amsmath}
\usepackage{amssymb}
\usepackage{pgfplots}
\pgfplotsset{compat=1.18}
\usetikzlibrary{calc,angles,quotes,arrows.meta,positioning,decorations.markings}
\begin{document}
\begin{tikzpicture}
\begin{axis}[
    width=13cm, height=9cm,
    xlabel={$x$}, ylabel={$y$},
    xmin=0, xmax=6, ymin=0, ymax=130,
    samples=200,
    axis lines=left,
    grid=major, grid style={gray!18},
    legend pos=north west,
    legend cell align=left,
    tick label style={font=\small},
    label style={font=\small},
]
\addplot[blue, very thick, domain=0:6] {x^3};
\addlegendentry{$y=x^3$（多项式）}
\addplot[red, very thick, domain=0:4.9] {exp(x)};
\addlegendentry{$y=e^x$（指数）}
% 交点附近：x^3=e^x 在 x≈4.536, y≈93.3
\addplot[only marks, mark=*, mark size=1.8pt, black] coordinates {(4.536,93.36)};
\node[anchor=south east, font=\small] at (axis cs:4.536,93.36) {交点 $x\approx4.54$};
% 注释框
\node[draw=red!60, fill=red!8, fill opacity=0.85, text opacity=1,
      rounded corners, align=left, font=\small, inner sep=4pt, anchor=east]
  at (axis cs:4.35,55)
  {交点之后 $e^x$ 一路碾压 $x^3$\\指数最终压倒任意多项式};
% 提示：交点前多项式暂时领先
\node[blue!60!black, anchor=west, font=\small] at (axis cs:0.45,40) {交点前 $x^3$ 暂时领先};
\end{axis}
\end{tikzpicture}
\end{document}
